\section{Requirement Alignment: Quorum vs. CometBFT}\label{sec:comparison}
This chapter evaluates how Quorum and CometBFT align with the requirements derived in Section~\ref{sec:requirements}. The focus is not on re-describing the SPS architecture or re-stating the requirements, but on understanding the practical match between those requirements and the design assumptions of the two platforms. In this sense, \emph{alignment} means that the platform's default mechanisms, operational model, and execution costs closely track the needs of a cyber-resilient SPS manager, without adding functionality that is either irrelevant or actively counterproductive.

The SPS manager is essentially a replicated, deterministic coordinator that receives sparse alerts, aggregates them, and triggers countermeasures. This is a narrow use case compared to the general-purpose decentralized application ecosystems that motivated many blockchain designs. Therefore, the most relevant dimension of comparison is whether each platform supports the required properties \emph{as first-class, low-overhead features}, and whether it avoids non-needed functionality that increases the trusted computing base or injects avoidable delay. Table~\ref{tab:comparison} summarizes the high-level mapping, the remainder of this section justifies each cell in that mapping and discusses the trade-offs that arise from different design philosophies.

\begin{table*}
\centering
\tblsetup
\caption{Mapping of Quorum and CometBFT features to SPS blockchain requirements.}
\label{tab:comparison}

\begin{tabularx}{\linewidth}{
  >{\raggedright\arraybackslash}p{4.1cm}
  >{\centering\arraybackslash}X
  >{\centering\arraybackslash}X
}
\toprule
\textbf{Property} & \textbf{Quorum (IBFT/QBFT)} & \textbf{CometBFT (ABCI)} \\
\midrule

\rowcolor{reqc}
\multicolumn{3}{l}{\textbf{Required}} \\
BFT consensus with finality     & \yes{Native}    & \yes{Native} \\
Permissioned static membership  & \yes{Native}    & \partialmark{Programmable} \\
Transaction signing             & \yes{Native}    & \yes{Native} \\
Event-driven block production   & \no{No}         & \yes{Configurable} \\[3pt]

\rowcolor{desc}
\multicolumn{3}{l}{\textbf{Desirable}} \\
In-memory replicated state      & \partialmark{Partial} & \yes{Flexible} \\
Low and predictable latency     & \partialmark{Fair}    & \yes{Excellent} \\
Graceful scaling with nodes     & \partialmark{Limited} & \partialmark{Programmable} \\
Configurable audit immutability & \partialmark{Limited} & \yes{Programmable} \\
Push-based notification         & \partialmark{Partial} & \yes{Programmable} \\[3pt]

\rowcolor{nac}
\multicolumn{3}{l}{\textbf{Not needed}} \\
Fixed-interval block production & \no{Present}    & \yes{Absent} \\
Gas / fee mechanisms            & \no{Present}    & \yes{Absent} \\
Token incentives                & \no{Inherited}  & \yes{Absent} \\
PoW / PoS                       & \no{Inherited}  & \yes{Absent} \\
Dynamic peer discovery          & \no{Present}    & \yes{Absent} \\
General-purpose VM              & \no{Present}    & \yes{Absent} \\
\bottomrule
\end{tabularx}
\end{table*}

\subsection{Alignment lens and scope}
The comparison is requirement-driven rather than feature-driven. It does not reward a platform for offering additional services, such as token economics or general-purpose smart contracts, unless those services directly improve the SPS manager's correctness, latency, or operational resilience. This is a deliberate lens: if a feature is not needed by the SPS manager, then its \emph{best} contribution is to be absent, because any additional subsystem increases configuration complexity and creates new failure modes that must be secured.

We also separate \emph{native} alignment (a property provided by the platform itself) from \emph{programmable} alignment (a property that can be implemented in the application layer). In SPS, programmable alignment can be acceptable, but only when the implementation effort is justified by measurable benefits such as lower latency, reduced attack surface, or greater determinism. This distinction is crucial when comparing a full-stack platform like Quorum with a consensus engine like CometBFT that intentionally pushes application logic to the developer.

\subsection{Required properties: hard constraints}

\subsubsection{BFT consensus with deterministic finality}
Quorum satisfies this requirement through IBFT and QBFT. Both protocols require a super-majority of validators (at least $2/3$) to sign each block, blocks are final and do not fork, and the system cannot progress if more than one third of validators are offline or Byzantine~\cite{QBFTCons90:online, moniz2020istanbul}. This behavior fits the SPS manager's safety needs because it provides a single authoritative state even under partial compromise.

CometBFT was explicitly designed to provide Byzantine fault-tolerant state machine replication with deterministic finality~\cite{buchman2018latest, cometbft}. It offers the same safety guarantees as IBFT/QBFT in terms of tolerating up to $f$ Byzantine validators in a set of $3f+1$. In practice, the key difference is not the safety model but where the application logic lives: Quorum merges consensus and application in the EVM, while CometBFT keeps consensus separate and exposes application logic through ABCI. For SPS, both satisfy the hard constraint, but CometBFT achieves it with less embedded application machinery.

\subsubsection{Permissioned static membership}
The SPS trust boundary is fixed, so validator membership should be permissioned and only change under explicit governance. Quorum supports this natively through validator set management in IBFT and QBFT. Validators are approved accounts, and membership can be changed through voting, either by JSON-RPC methods or via a governance smart contract~\cite{QBFTCons90:online}. This closely matches the requirement, but the contract-based approach introduces additional on-chain logic that must itself be secured.

CometBFT does not include a built-in permissioning layer. Membership is defined in the genesis configuration and can be updated by the application through ABCI. This means static membership can be enforced as a policy in the application logic, which is adequate for SPS but requires additional engineering. The trade-off is clear: Quorum delivers the feature out of the box, while CometBFT makes the designer explicitly encode the governance rules that fit the SPS deployment. In a research-grade SPS, this explicitness can be advantageous because it avoids coupling permissioning to generic smart-contract governance.

\subsubsection{Transaction signing and accountability}
Transaction signing is a native property of the Ethereum lineage on which Quorum is built, every transaction carries a signature and is authenticated by the consensus layer. For the SPS manager, this provides accountability for IDS nodes and actuators, enabling non-repudiation and forensic auditing.

In CometBFT, transaction validation is explicitly delegated to the application through ABCI. The standard pipeline calls \texttt{CheckTx} for mempool admission and later \texttt{ProcessProposal} or \texttt{FinalizeBlock} for deterministic execution, and these hooks are the natural place to verify signatures and enforce authorization policies~\cite{abciOverview}. This yields the same functional property as in Quorum, but the responsibility for signature schemes and key management is pushed to the application. That extra work is acceptable in SPS because signing policies are tightly linked to the SPS trust model rather than to generic blockchain account models.

\subsubsection{Event-driven block production}
Alert traffic in SPS is sparse and bursty, which makes fixed-interval block production a poor fit. Quorum's QBFT configuration uses a block period timer (\texttt{blockperiodseconds}) that triggers proposal attempts at a minimum fixed interval~\cite{QBFTCons90:online}. This keeps consensus active even under idle conditions and introduces a baseline delay between detection and commitment that is independent of application complexity.

CometBFT allows disabling empty block creation through configuration, effectively enabling blocks only when transactions exist~\cite{cometbftConfig}. This matches the SPS requirement for event-driven progression: the system can remain quiescent without consuming consensus cycles, and when an alert arrives it can be committed without waiting for the next periodic block. In requirement alignment terms, CometBFT is closer to the desired behavior and can be tuned to minimize idle overhead.

\subsection{Desirable properties: efficiency and responsiveness}

\subsubsection{Execution model and deterministic state transitions}
Quorum inherits the Ethereum execution model. Transactions are RLP-encoded, executed in the EVM, and charged gas at instruction granularity~\cite{buterin2013ethereum}. Even in permissioned deployments where fees are irrelevant, the EVM still enforces gas accounting and executes bytecode in a general-purpose virtual machine. The SPS response logic, however, is a small deterministic automaton, it uses only a narrow subset of EVM expressiveness while paying the full runtime cost and complexity of a general-purpose environment.

CometBFT externalizes execution through ABCI, allowing the SPS state machine to be implemented in a native language such as Go or Rust. This eliminates bytecode interpretation and gas accounting overhead, and it makes the runtime cost of a state transition proportional to its intrinsic logic rather than to virtual machine scaffolding. ABCI also imposes strict determinism requirements on several handlers (for example, \texttt{ProcessProposal} must be deterministic)~\cite{abciOverview}. This constraint is compatible with SPS, which already requires deterministic logic, but it places additional discipline on the implementation to avoid non-deterministic libraries, wall-clock time, or concurrency races that would diverge across replicas.

\subsubsection{Transaction and state modeling}
SPS alerts and response decisions are naturally expressed as structured, domain-specific messages: they carry identifiers for monitored variables, confidence scores, sensor provenance, and explicit transition requests. Quorum forces these messages into the Ethereum transaction model, where payloads are encoded as calldata and interpreted by a contract. This works, but it introduces a semantic gap: the meaning of an alert is not part of the transaction itself, but is reconstructed by contract code that parses byte arrays and manages storage slots. The resulting system is functional, yet the data model is indirect and tends to be more verbose than the conceptual state-machine representation.

CometBFT's transaction model is application-defined, which means the SPS designer can encode alerts as explicit data structures, perform validation in \texttt{CheckTx}, and map them directly to state transitions in \texttt{FinalizeBlock}~\cite{abciOverview}. This reduces parsing overhead and preserves semantic clarity throughout the pipeline. It also enables the inclusion of structured metadata that supports assurance claims, such as explicit evidence tags or signed measurement bundles, without forcing them through EVM-specific encoding constraints. In terms of requirement alignment, the ability to define the transaction schema is an advantage because it lets the blockchain layer reflect the SPS domain model rather than the other way around.

\subsubsection{State storage and in-memory replication}
The SPS state is small and stable, which makes in-memory replication attractive. Quorum, as an Ethereum client, persists its state in a database optimized for large smart-contract state and historical storage. This persistence is beneficial for general-purpose applications, but in SPS it adds I/O overhead and enlarges the system's storage footprint without clear benefit, especially when the response logic only needs the current state and short-term history.

CometBFT separates consensus from application state, so the application is free to choose storage strategies. The core node still persists its consensus metadata and block store, but the application can keep its own replicated state in memory and persist only what is operationally required. This flexibility aligns with the desirable requirement for in-memory state and opens the possibility of tailoring retention policies to resource-constrained deployments.

\subsubsection{Latency predictability at low load}
SPS effectiveness depends on short and predictable response time. In Quorum, two mechanisms contribute to baseline latency: the fixed block period and the EVM execution overhead. Even when only a single alert transaction is submitted, the transaction must wait for the next block proposal window, consensus then proceeds, and the state transition executes in the EVM with gas metering. This yields stable but non-minimal latency that scales with configured block time rather than with alert rate.

In CometBFT, latency is dominated by the consensus round and the application execution time. When empty blocks are disabled, the first transaction can be proposed immediately without waiting for a periodic timer. For a small validator set, the BFT round completes quickly, and the native application logic executes with minimal overhead. This matches the SPS preference for predictable, low-latency decisions, especially under sparse alert traffic.

\subsubsection{Burst handling and backpressure}
SPS workloads are not only sparse but also bursty: coordinated attacks can trigger multiple IDS alerts in a short window, and a robust system must absorb these bursts without collapsing into unbounded latency. In Quorum, the fixed block period provides a stable scheduling rhythm, but it also places a hard cap on how quickly the system can advance the ledger. When alerts arrive faster than blocks can be proposed, transactions accumulate in the pool and are executed in subsequent blocks, increasing end-to-end delay.

CometBFT exposes more opportunities for backpressure and filtering at the application boundary. Because \texttt{CheckTx} is invoked before a transaction enters the mempool, an SPS application can reject duplicates, apply rate limits, or enforce priority rules based on IDS provenance and severity~\cite{abciOverview}. This capability does not eliminate the fundamental consensus bottleneck, but it allows the designer to manage bursts in a way that preserves the most critical alerts. Such application-level control is aligned with the SPS requirement for bounded latency under attack conditions, even though it requires explicit policy design.

\subsubsection{Scaling with validator count}
BFT consensus protocols exhibit $O(n^2)$ communication in the voting phase, which limits scalability as the validator set grows~\cite{wang2022bft}. Quorum's IBFT/QBFT adheres to this pattern: validators exchange prepare and commit messages with all others, and the protocol stalls if more than one third of validators are offline. This is acceptable for small SPS deployments but constrains scale.

The scalability tension is not incidental, it is part of a broader trade-off among decentralization, security, and performance that appears in many blockchain designs~\cite{monte2020scaling}. SPS deployments typically prioritize security and availability over large-scale decentralization, which makes small validator sets acceptable. However, even within this constrained regime, the system should degrade gracefully as additional validators are added for resilience. Quorum's protocol does not change its message complexity as the validator set grows, so communication cost increases quadratically with each added node.

CometBFT uses gossip dissemination to broadcast proposals and transactions, which reduces the leader's broadcast burden and can lower propagation latency compared to direct full-mesh dissemination~\cite{buchman2018latest}. However, the final voting phase still requires super-majority signatures, and the fundamental $O(n^2)$ message pattern remains. In requirement terms, neither platform offers a built-in path to graceful scaling with large validator sets. CometBFT's advantage is that it makes validator-set management programmable, enabling committee rotation or hierarchical validation to be implemented at the application layer if such strategies are desired.

\subsubsection{Configurable audit immutability}
Auditability is useful for forensic analysis, but in SPS it is not always necessary to retain complete transaction history indefinitely. Quorum, following Ethereum's model, is optimized for immutable ledger history and full archival state, which is consistent with public-blockchain assumptions but potentially excessive for closed, resource-constrained SPS deployments.

CometBFT's separation of concerns enables a more flexible policy: the application can decide which events must be recorded, which can be summarized, and which can be discarded after a retention window. This does not weaken consensus safety, but it does require explicit design to ensure that the retained data still supports accountability and audit obligations. From an alignment perspective, CometBFT better supports configurable immutability because it does not assume a fixed archival model.

\subsubsection{Push-based notifications}
SPS actuators benefit from push-based notification, since polling introduces additional latency. Quorum supports event logs and client subscriptions through the Ethereum API, which can be used to trigger actuation when the smart contract emits a state-change event. This capability is mature and widely supported in Ethereum tooling.

In CometBFT, events can be emitted by the application during block finalization and consumed by clients through the node's RPC event subscription mechanism. While this requires explicit wiring in the application layer, it gives the SPS designer direct control over the semantics and granularity of notifications. Both platforms can support push-based actuation, but CometBFT again requires application-level implementation rather than offering it as a fixed platform feature.

\subsection{Alignment with the SPS control loop}
Self-protecting systems are commonly framed as a control loop (e.g., Monitor, Analyze, Plan, Execute) operating on a shared knowledge base~\cite{kephart2003vision, yuan2014systematic}. The blockchain layer, in this context, is not a general ledger but the coordination substrate of that loop. Alignment therefore depends on how directly the platform supports the control loop's timing and determinism. Event-driven consensus is a closer match to the monitor--analyze cadence of SPS than fixed block periods, because it avoids periodic consensus activity when no security-relevant events exist.

The SPS manager's internal state is a compact representation of monitored variables and response states, which can be modeled as a finite-state machine~\cite{iannucci2018model}. Quorum allows this logic to be expressed in Solidity, however, the EVM representation tends to be more verbose than the conceptual model, and the state machine structure is often obscured by contract bookkeeping, gas accounting, and event emission. CometBFT permits the same state machine to be implemented as a direct data structure with explicit transitions, which makes the state logic more transparent and easier to map to the formal model used during SPS design.

The MAPE-K perspective also highlights how knowledge updates should be triggered. In an SPS, the knowledge base should reflect the most recent corroborated evidence, not the arrival of periodic timing events. Quorum's periodic block production effectively introduces a time-based update cadence even when no new evidence exists, whereas CometBFT can be configured to update the shared state only when new evidence appears. This difference does not change correctness, but it affects how tightly the blockchain layer is coupled to the semantics of monitoring and analysis, and therefore how cleanly the SPS control loop can be interpreted.

The voting and corroboration rules that protect against compromised IDSs can be placed in different phases of the transaction pipeline. In Quorum, the typical implementation encodes those rules inside the smart contract, so every alert travels through consensus before being filtered. This preserves a complete audit trail, but it also increases the volume of consensus traffic and imposes EVM execution overhead for inputs that will ultimately be discarded. In CometBFT, the ABCI interface allows preliminary checks in \texttt{CheckTx} and proposal-level validation in \texttt{ProcessProposal} before finalization~\cite{abciOverview}. This enables the designer to reject redundant or malformed alerts early while still preserving deterministic behavior, which aligns with the requirement for low-latency actuation in the presence of noisy sensors.

From an assurance perspective, the SPS community emphasizes the need for evidence that adaptation logic is correct and remains correct at runtime~\cite{calinescu2017engineering, pekaric2023systematic}. A concise, deterministic state machine makes it easier to construct such evidence and to reason about the impact of each transition. Quorum's smart-contract deployment supports auditability but makes it harder to separate the state machine itself from the execution environment. CometBFT's explicit separation of consensus and application code makes that boundary visible, which supports assurance arguments about the correctness of the control logic and its traceability to the SPS design artifacts.

\subsection{Non-required features and attack surface}
Several features inherited by Quorum are classified as not needed in Section~\ref{sec:requirements}. Most notably, Quorum retains the EVM and gas accounting, which are essential for public Ethereum but unnecessary for a permissioned SPS coordinator. These features do not merely sit idle, they increase runtime overhead and enlarge the trusted computing base. They also expose the SPS logic to known smart-contract vulnerability classes, such as reentrancy, integer overflow patterns, and access-control errors documented in surveys of Ethereum contract vulnerabilities~\cite{atzei2017survey, kushwaha2022systematic}.

Quorum additionally provides a privacy subsystem based on Tessera, along with an enclave-based cryptographic separation for private transaction data~\cite{goquorumPrivacy}. These components are powerful for enterprise privacy use cases, but in SPS they are largely orthogonal: IDS alerts and actuation decisions are intended to be shared among all validators, not hidden. Introducing a private-transaction manager adds extra services, configuration, and key-management procedures that do not contribute to the SPS goals and therefore reduce requirement alignment.

CometBFT avoids these inherited features by design. It provides consensus, networking, and a minimal transaction processing pipeline, and it leaves economic incentives, virtual machines, and privacy overlays to the application. This absence of non-required functionality is not a limitation in the SPS context, it is a form of alignment, because it keeps the coordination layer narrow and reduces the amount of code that must be trusted for correct security decisions.

\subsection{Visibility, confidentiality, and the trust boundary}
The SPS manager exists to produce a shared, authoritative view of monitored security state, which implies that validators must see the same alerts and agree on the same response decisions. This pushes the design toward transparency rather than confidentiality: hiding alerts from some validators would weaken the evidence base and complicate consensus. In this sense, a fully shared state is aligned with the SPS trust boundary, and strong privacy mechanisms are not generally required.

However, some deployments may treat IDS alerts as sensitive because they reveal vulnerabilities or operational details about the managed system. Quorum's privacy architecture, based on Tessera and enclave-backed cryptographic separation, provides a ready-made mechanism for restricting transaction visibility to selected participants~\cite{goquorumPrivacy}. This can be attractive in regulated environments, but it also changes the semantics of shared state: validators that cannot see private transactions do not execute them, which can lead to divergent local views unless additional coordination is implemented. For SPS, where shared awareness is central, such divergence must be carefully managed, and the privacy layer becomes a significant part of the trusted computing base.

Privacy also interacts with accountability. If alerts are visible only to a subset of validators, then audit trails become fragmented and attribution becomes harder to verify across the full set of replicas. This undermines one of the implicit advantages of using a blockchain in SPS: a single, consensus-backed record of security events. The design therefore faces a trade-off between confidentiality and collective verifiability, and the chosen platform should make that trade-off explicit rather than implicit. Quorum's privacy extensions are powerful but can obscure this tension if adopted without a clear policy.

CometBFT does not include a built-in privacy layer, but its application-defined transaction model allows confidentiality mechanisms to be designed around the SPS needs. For example, alerts can be encrypted while still carrying a cleartext envelope that indicates category or severity, allowing validators to reach consensus on high-level state while protecting details. This approach requires cryptographic design and policy decisions, yet it keeps the confidentiality model explicit and aligned with the system's governance. The key point for requirement alignment is that confidentiality is not free: whichever platform is chosen, it must be reconciled with the need for a coherent shared state, and CometBFT makes that reconciliation a deliberate design choice rather than an inherited default.

\subsection{Security and assurance implications}
The SPS manager is a critical component: errors in its logic can lead to false positives, missed attacks, or unsafe countermeasures. When implemented as a smart contract in Quorum, the response logic inherits the risk profile of Ethereum contracts. Extensive analyses show that subtle bugs in contract logic are common and can lead to serious vulnerabilities~\cite{atzei2017survey, kushwaha2022systematic}. Formal verification and defensive coding can mitigate this, but they require specialized expertise and still operate within the constraints of the EVM model.

The self-adaptive systems literature emphasizes that correctness arguments should not end at design time, they must be maintained at runtime through explicit assurance mechanisms~\cite{calinescu2017engineering, pekaric2023systematic}. In this view, the SPS coordinator is not simply a data plane but an assurance-critical component that should produce evidence about its own decisions. Quorum's smart-contract model can expose such evidence through immutable logs, yet the opacity of the EVM and the complexity of contract execution can make it difficult to relate that evidence to the high-level SPS assurance arguments. A verifier must reason not only about the intended state machine but also about EVM execution semantics, gas limits, and compiler effects.

CometBFT enables the same logic to be implemented as native code, where conventional software engineering techniques, static analysis tools, and language-level safety guarantees can be applied. This does not eliminate risk, but it avoids EVM-specific pitfalls and makes it easier to integrate the SPS logic with existing, well-tested security libraries. The key assurance challenge shifts from smart-contract correctness to deterministic replication: application code must be deterministic across replicas, or else consensus will diverge. ABCI makes this requirement explicit, which aligns with the deterministic nature of SPS response logic but must be carefully enforced in the implementation~\cite{abciOverview}.

Validator governance introduces another assurance dimension. Quorum's on-chain validator management means that governance decisions are themselves encoded in contracts, which become part of the trusted base. If the governance contract is flawed or if validator key management is weak, the system's trust boundary can be altered by an attacker. In CometBFT, validator governance is an application policy, this allows integration with existing organizational controls (for example, manual enrollment, hardware-backed keys, or external approval workflows), but it also requires a careful design to ensure that membership changes are auditable and resistant to compromise.

\subsection{Failure modes, recovery, and liveness}
Both Quorum and CometBFT are BFT protocols and therefore inherit the classical liveness constraint: if more than one third of validators are Byzantine or offline, consensus stalls. The GoQuorum documentation makes this explicit for QBFT/IBFT and highlights that block production stops when the threshold is exceeded~\cite{QBFTCons90:online}. In an SPS deployment, such stalls are not merely availability issues, they directly delay countermeasures, which can degrade the system's defensive posture. This makes validator health monitoring and rapid recovery procedures part of the requirement alignment story, not just an operational afterthought.

Quorum's time-driven block production can also mask the difference between a healthy-but-idle system and a partitioned system, because the consensus layer continues to propose blocks on a schedule. This behavior is not incorrect, but it generates activity even when no alerts exist and can complicate the interpretation of liveness signals. CometBFT's ability to disable empty blocks~\cite{cometbftConfig} is more compatible with an alert-driven workload because consensus activity becomes a clearer indicator of meaningful events, and idle periods are expected rather than anomalous.

Timeout behavior further differentiates alignment with SPS latency requirements. In QBFT/IBFT, the protocol advances through rounds with a configurable \texttt{requesttimeoutseconds} value and doubles that timeout on consecutive failures, which increases reaction time under partial failure~\cite{QBFTCons90:online}. This backoff is appropriate for public or unstable networks, but in a tightly controlled SPS deployment it can create long delays in the very situations where fast countermeasures are most needed. CometBFT follows a similar round-based structure but exposes more control at the application and configuration level, allowing the designer to tune timeouts to the expected alert profile and network conditions. This tuning is part of requirement alignment because it ties consensus responsiveness to the SPS threat model rather than to generic blockchain defaults.

Recovery behavior depends on how quickly a node can rejoin and reconstruct the shared state. Quorum's full-history model is conservative and strongly consistent, but it also tends to require more data transfer to resynchronize a new or recovering validator. CometBFT allows the application to define snapshotting or compact state representations that can accelerate recovery, at the cost of additional design complexity. This trade-off mirrors the overall alignment pattern: Quorum offers stronger defaults but heavier operational cost, while CometBFT offers flexibility that must be explicitly engineered.

From the perspective of trustworthy autonomous systems, resilience is a composite property that includes not only correctness during normal operation but also predictable recovery under failure~\cite{flammini2022towards}. The requirement alignment question therefore includes how transparent and controllable the recovery process is. CometBFT's explicit application boundaries make it easier to reason about which components must recover and how state is reconstructed, whereas Quorum's monolithic stack can obscure these boundaries behind the EVM and database layers.

\subsection{Operational fit and lifecycle considerations}
From an operational perspective, Quorum is a full-stack platform with mature tooling, JSON-RPC APIs, and a wide ecosystem inherited from Ethereum. This reduces implementation effort for generic blockchain applications but can be excessive for SPS. The SPS use case is narrow, and the operational burden of maintaining additional components (privacy managers, contract deployment infrastructure, smart-contract upgrade procedures) may outweigh their benefits.

CometBFT requires more engineering effort up front because the application layer must be built explicitly. However, this effort directly targets the SPS needs: the designer can implement only the required transaction formats, validation logic, and state transitions, without accommodating generic smart-contract execution. The resulting system is smaller and easier to reason about, which is a form of operational alignment even if it increases initial development effort.

Upgrade paths also differ. Quorum can evolve through contract upgrades or proxy patterns, which allow in-place logic changes but introduce governance complexity and risk of misconfiguration. CometBFT applications typically evolve through coordinated binary upgrades, which are operationally heavier but explicit. For SPS, explicit upgrade coordination may be preferable because it ensures that all validators run the same deterministic code and that updates are reviewed with the same rigor as other security-critical changes.

Integration with IDSs and actuators also highlights the difference between a general-purpose platform and a purpose-built state machine. Quorum's interfaces are standardized and compatible with existing Ethereum tooling, which can simplify client development, but they expose a large API surface that the SPS components must understand and secure. A CometBFT-based SPS can expose narrow, domain-specific APIs that mirror the alert schema and actuation semantics, reducing the potential for misuse and simplifying validation at the system boundary. This is not a technical superiority on its own, but it aligns with the principle of least privilege that underlies secure autonomic systems.

Tooling maturity is not irrelevant to requirement alignment. Quorum benefits from extensive Ethereum tooling for contract testing, security auditing, and monitoring, which can accelerate early prototypes. The downside is that this tooling assumes an EVM-centric world and encourages design patterns (upgradable contracts, modular governance tokens, complex on-chain registries) that are misaligned with the minimal SPS coordinator model. CometBFT instead leverages conventional software engineering toolchains. This shifts the burden to application developers but also gives them the ability to adopt mature testing and verification frameworks from mainstream languages without translating logic into smart contracts.

Operational observability is another alignment factor. In SPS, the blockchain layer is part of the security apparatus, so logs and metrics must support incident analysis and forensic reconstruction. Quorum exposes Ethereum-style logs and tracing, which provide a rich event stream but at the cost of higher volume and lower semantic specificity. A CometBFT-based SPS can emit domain-specific events with explicit semantics, which can simplify forensic reasoning and support the dynamic assurance arguments advocated in the self-adaptive systems literature~\cite{calinescu2017engineering}. This flexibility again shifts responsibility to the application designer but allows the observability layer to be aligned with the system's threat model.

\subsection{Synthesis: alignment profile}
Both platforms satisfy the core safety requirement of Byzantine fault tolerance and finality. The difference lies in how much non-essential machinery each platform brings along. Quorum offers a rich execution environment and enterprise features, but those features are largely unneeded in SPS and therefore impose avoidable overhead and risk. Its fixed block period and EVM-centric execution model are mismatched with the sparse, deterministic, low-latency profile of SPS coordination.

CometBFT aligns more closely with the requirement set because it provides consensus and leaves the rest to the application designer. This yields a system that can be tailored to the SPS state machine, can avoid empty blocks, and can keep the trusted computing base small. The cost of this alignment is the need to engineer application-level policies for permissioning, notification, and audit retention. In a research or prototype SPS, this trade-off is acceptable, and it explains why CometBFT emerges as the more requirement-aligned technology despite offering fewer out-of-the-box features.

The comparison does not imply that Quorum is categorically unsuitable. If a deployment values compatibility with Ethereum tooling, requires smart-contract composability, or must leverage privacy features such as Tessera, Quorum can provide practical benefits that offset its misalignment with SPS-specific requirements. Conversely, if the primary goal is a minimal, deterministic coordinator with tight control over execution and latency, the additional engineering effort of CometBFT is justified. Requirement alignment therefore becomes a choice about which costs are acceptable: operational complexity and extra subsystems on the Quorum side, or increased development responsibility on the CometBFT side. This framing is useful for thesis-level evaluation because it clarifies that the platform decision is fundamentally about matching system assumptions, not simply about raw performance.
